\chapter{拓扑空间简介}

\begin{tikzpicture}[align=center]
    \node[] (1) at(0,2) {集合};
    \node[] (2) at(3,2) {\color{blue} 集合间的关系：映射};
    \node[] (3) at(2,1) {距离};
    \node[] (4) at(5/2,0) {度量空间};
    \node[] (5) at(4,-1) {无法定义距离：给出各种程度的“附近”};
    \node[] (6) at(0,-2) {拓扑空间};
    \node[] (7) at(3,-2) {\color{blue} 连续性};
    \node[] (8) at(1,-3) {微分结构};
    \node[] (9) at(0,-4) {微分流形};
    \node[] (10) at(4,-4) {\color{blue} 导函数的连续性};
    \draw [->] (1) -- (2);
    \draw [->] (1) -- (4);
    \draw [->] (1) -- (6);
    \draw [->] (6) -- (7);
    \draw [->] (6) -- (9);
    \draw [->] (9) -- (10);
\end{tikzpicture}

\section{集论初步}

\subsection{集合}

\begin{itemize}
    \item 集合，元素，子集
    \item 集合的表示
    \begin{itemize}
        \item 列举法
        \item 描述法
    \end{itemize}
    \item 集合的运算：交，并，差，补
    \begin{itemize}
        \item 运算律
    \end{itemize}
    \item Cartesian积
    \item Euler度量
\end{itemize}

\subsection{映射}

\begin{itemize}
    \item 映射的分类
    \begin{itemize}
        \item 单射
        \item 满射
        \item 双射
        \item[$\cdot$] 常值映射，复合映射
    \end{itemize}
    \item 连续映射：将子空间拓扑下的开集映射到开集
\end{itemize}

\section{拓扑空间}

\begin{itemize}
    \item 拓扑，拓扑空间，开集，闭集
    \begin{itemize}
        \item 离散拓扑
        \item 凝聚拓扑
        \item 通常拓扑
        \item 乘积拓扑
        \item 诱导拓扑，拓扑子空间
    \end{itemize}
    \item 连续性
    \item 同胚
    \item 元素的邻域，子集的邻域，开邻域，内点
    \begin{itemize}
        \item[\o] 开集$\iff$每一点都是内点
    \end{itemize}
    \item 连通性
    \item 闭包，内部，边界
\end{itemize}

\section{紧致性}

\begin{itemize}
    \item 覆盖，开覆盖，有限子覆盖
    \item 紧致：任意开覆盖都有有限子覆盖
    \begin{itemize}
        \item[\o] 连续映射保持紧致性
    \end{itemize}
    \item $T_2$拓扑空间
    \begin{itemize}
        \item[\o] $T_2$空间，紧致子集$\implies$闭集
        \item[\o] 紧致拓扑空间，闭子集$\implies$紧致
    \end{itemize}
    \item 有界性
    \item 拓扑性质：紧致性，连通性，$T_2$性质
    \begin{itemize}
        \item[\o] 有界性定理，最值定理
        \item $\mathbb{R}^n$上紧致子集$\iff$有界闭集
    \end{itemize}
    \item 序列，极限，聚点，第二可数
    \begin{itemize}
        \item 紧致子集$\implies$任意序列都有包含在内的聚点\\
              第二可数集合的子集，若任意序列都有包含在内的聚点$\implies$紧致
    \end{itemize}
\end{itemize}